% Generated draft: requires mathematical and source review.
\documentclass[11pt]{article}
\usepackage[margin=25mm]{geometry}
\usepackage{fontspec}
\setmainfont{Latin Modern Roman}
\usepackage{amsmath,amssymb,mathtools,unicode-math}
\setmathfont{Latin Modern Math}
\usepackage{enumitem,needspace,xurl,hyperref}
\setlength{\parindent}{0pt}
\setlength{\parskip}{5pt}
\setlength{\emergencystretch}{3em}
\newcommand{\N}{\mathbb{N}}
\newcommand{\Z}{\mathbb{Z}}
\newcommand{\Q}{\mathbb{Q}}
\newcommand{\R}{\mathbb{R}}
\newcommand{\C}{\mathbb{C}}
\newcommand{\F}{\mathbb{F}}
\newcommand{\A}{\mathbb{A}}
\newcommand{\PP}{\mathbb{P}}
\newcommand{\E}{\mathbb{E}}
\newcommand{\eps}{\varepsilon}
\newcommand{\dd}{\,\mathrm{d}}
\begin{document}
\section*{UnsolvedMath ID: 30002943}
\textbf{Source problem code:} catalogue entry
\subsection{Definitions and mathematical statement}
Let \(B_1\) and \(B_2\) be Banach spaces, and let \(T_N:B_1\to B_2\), \(N\in\N=\{1,2,\ldots\}\), be bounded linear operators. The scalar field is not specified in the supplied record. Let \(D\subseteq B_1\) be a dense subset, meaning that \(\overline D^{\|\cdot\|_{B_1}}=B_1\). Boundedness is stated for each operator separately; no uniform bound on \(\|T_N\|\) is supplied.

The record says that the operators converge on a dense subset, but it does not define the target \(T\) appearing in the requested conclusion, specify its properties, or explicitly identify the convergence on the dense subset with convergence to \(T\). Consequently, the following is a conditional formulation of the intended question, not a completed specification of its hypotheses.

If a target map \(T:B_1\to B_2\) is specified and the dense-subset convergence is intended to mean
\[
\lim_{N\to\infty}\|T_Ng-Tg\|_{B_2}=0
\qquad\text{for every }g\in D,
\]
does every \(f\in B_1\) admit a strictly increasing sequence of positive integers \((N_k(f))_{k\geq1}\) such that
\[
\lim_{k\to\infty}\|T_{N_k(f)}f-Tf\|_{B_2}=0?
\]
More explicitly, the requested subsequence property has quantifiers
\[
\forall f\in B_1\quad
\exists (N_k(f))_{k\geq1}\subseteq\N\quad
\left[N_1(f)<N_2(f)<\cdots\ \text{ and }\
\forall\varepsilon>0\ \exists k_0\in\N\ \forall k\geq k_0:
\|T_{N_k(f)}f-Tf\|_{B_2}<\varepsilon\right].
\]
The subsequence may depend on \(f\); the record does not request a single subsequence valid for all \(f\). It allows divergence away from the dense subset but does not require divergence to occur. It remains unclear whether the question concerns every sequence satisfying the stated abstract hypotheses or a particular class of approximation schemes from the cited workshop discussion. No additional hypotheses defining such a class are supplied, and none are imposed here.
\subsection{Short English statement}
For each input, can one choose its own subsequence of approximation operators whose outputs converge in norm to the intended target value? The supplied formulation leaves the target operator and possibly the intended class of approximation schemes unspecified, so this is a conditional statement requiring mathematical review.
\subsection{Sources}
\begin{itemize}
\item Local registry: \texttt{data/unsolvedmath/problems.json}, supplied record with permanent UnsolvedMath ID \(30002943\), title ``Pointwise Subsequence Convergence of Approximation Operators,'' source problem code \texttt{catalogue entry}, in the Oberwolfach Reports Open Problems collection. No individual proposer is named; the supplied proposed year is \(2015\).
\item Supplied source citation: \emph{Applied Harmonic Analysis and Sparse Approximation} (2016), Oberwolfach Reports, pp.~2189--2263, DOI \texttt{10.4171/owr/2015/38}. \url{https://doi.org/10.4171/owr/2015/38}. The background attributes the problem to workshop open-problem discussions documented in Oberwolfach Reports. No problem-specific page or passage is supplied.
\item Supplied literature source, labelled ``Boche and Pohl 2016'': \url{https://www.sciencedirect.com/science/article/pii/S0021904516000034}. The record supplies no article title or pinpoint reference.
\item Supplied literature source, labelled ``original OWR report'': \url{https://doi.org/10.4171/owr/2015/38}.
\item Literature assessment dated exactly \texttt{2026-08-21}: the supplied curation labels the problem open and reports strong divergence and finite search horizon obstructions for important approximation schemes, while describing the universal function dependent subsequence question as unresolved. This is supplied dated curation, not an independently verified literature assessment. The background explicitly calls for independent verification of time-sensitive status and cited claims.
\item The registry also supplies the statement status \texttt{corrected\_verified}, with an assertion that mathematical content was preserved while notation, typography, or omitted definitions were repaired. This is source metadata, not independent mathematical verification of the present draft.
\end{itemize}
\end{document}
